%%% historicity-baseline.tex — Bayesian baseline supplement
%%% Compile: xelatex historicity-baseline && xelatex historicity-baseline
%%%
%%% Corpus-linguistic baseline for adelphos in contemporary-or-earlier
%%% non-Christian Greek, for the Bayesian prior argument of the main
%%% paper (§3.5).

\documentclass[11pt,oneside]{article}

\usepackage[letterpaper,margin=1in,headheight=14pt]{geometry}
\usepackage{fontspec}
\usepackage{xcolor}
\usepackage{fancyhdr}
\usepackage{booktabs}
\usepackage{enumitem}
\usepackage{array}
\IfFileExists{csquotes.sty}%
  {\usepackage[autostyle,english=american]{csquotes}}%
  {\providecommand{\enquote}[1]{``##1''}}
\usepackage[colorlinks=true,linkcolor=blue!60!black,urlcolor=blue!60!black]{hyperref}

\defaultfontfeatures{Ligatures=TeX}
\setmainfont{STIX Two Text}

\pagestyle{fancy}
\fancyhf{}
\fancyhead[L]{\small\itshape Baseline: \textit{adelphos} in contemporary Greek}
\fancyhead[R]{\small\thepage}
\renewcommand{\headrulewidth}{0.4pt}

\setlength{\parindent}{0pt}
\setlength{\parskip}{0.4em}

\begin{document}

\begin{center}
  {\LARGE\bfseries Supplementary Material:\\[4pt]
   Corpus-Linguistic Baseline for \textit{Adelphos}\\[2pt]
   in Contemporary Greek}\\[18pt]
  {\large Companion to \textit{The Evidentiary Status of the Historical Jesus}}\\[6pt]
  {\normalsize Stan Thomas}\\[4pt]
  {\normalsize Independent Researcher, St.\ George, Utah\\
   ORCID: 0000-0002-8828-7856}\\[12pt]
  {\normalsize Draft of \today}
\end{center}

\vspace{0.5em}

\section*{Purpose}

The main paper's §3.5 observes that across the seven undisputed Pauline
letters, 117 of 121 \textit{adelph-} tokens classify as community/fellow-
believer, zero as unambiguously biological-sibling, and four as disputed
(Rom 9:3, Rom 16:15, 1 Cor 9:5, Gal 1:19).  That observation takes on
evidential force only if we can say how \textit{atypical} Paul's
distribution is against a background of contemporary-or-earlier Greek
usage, uncontaminated by any influence from Paul himself or from later
Christian discourse.

This supplement establishes that background.  It samples
\textit{adelph-} family occurrences across six non-Christian corpora
spanning pre-Pauline to contemporary non-Christian Greek, classifies
each token into the same A/B/C/D/E scheme used in the companion
\textit{Pauline Occurrences} supplement (extended with categories D and
E for ethnic kinship and literary metaphor), and reports per-bucket and
aggregate proportions with Wilson 95\% confidence intervals.

\section*{Corpus design}

Six \textbf{non-Christian} baseline buckets, chosen to cover the
register diversity Paul would have encountered (documentary,
scriptural, historiographical, philosophical, literary) and to include
both pre-Pauline and contemporary non-Christian sources.  Two
additional \textbf{post-Pauline Christian} buckets are included for
contextual comparison \textit{only}---they are downstream of Pauline
discourse and therefore cannot inform a prior.

\begin{center}
\begin{tabular}{@{}lll@{}}
\toprule
Bucket & Dates & Rationale \\
\midrule
\multicolumn{3}{l}{\textit{\textbf{Baseline — non-Christian, contemporary or earlier}}} \\
Papyri (documentary) & 100 BCE – 70 CE & Everyday Greek; expected heavy B \\
Septuagint (LXX) & 3rd–1st c.\ BCE & Paul's scripture; heavy B and D \\
Philo of Alexandria & c.\ 20 BCE – 50 CE & Hellenistic Jewish-Greek, contemporary \\
Josephus & writing 75–95 CE & Jewish, non-Christian, contemporary \\
Pre-Pauline literary & 2nd c.\ BCE – 1st c.\ BCE & Polybius, Diodorus, Strabo (pagan) \\
Hellenistic literary & c.\ 50–120 CE & Plutarch, Epictetus, Dio Chrysostom (pagan) \\
\midrule
\multicolumn{3}{l}{\textit{Downstream Christian — for context, not prior}} \\
NT non-Pauline letters & c.\ 60–100 CE & Hebrews, James, 1–2 Peter, 1–3 John, Jude \\
Gospels + Acts & c.\ 70–95 CE & Narrative Christian; biographical B-references \\
\bottomrule
\end{tabular}
\end{center}

Each bucket contributes $n \approx 35$ randomly-sampled tokens, for a
baseline total of $N = 212$.  Sampling is by
\texttt{random.sample}(seed=42) over the full set of \textit{adelph-}
family occurrences identified per corpus.

Sources: SBL Greek New Testament; Rahlfs Septuagint; Loeb/Niese
Josephus (via Perseus); Duke Databank of Documentary Papyri (via GitHub
EpiDoc repository); Cohn--Wendland Philo (via \textit{Open Greek and
Latin}); Perseus canonical-greekLit for literary authors.  All sources
are openly licensed or public-domain.

\section*{Classification scheme}

Identical to the Pauline supplement, extended with D (ethnic/tribal/
covenantal kinship) and E (literary/metaphorical extension) because the
broader corpus routinely exhibits these uses where Paul almost never
does.

\begin{description}[leftmargin=2em]
\item[\textbf{A}] Community / voluntary association (religious,
  philosophical, professional, cultic).  Membership by election.
\item[\textbf{B}] Biological sibling (shared parentage).
\item[\textbf{C}] Disputed / ambiguous (borderline cases default here).
\item[\textbf{D}] Ethnic / tribal / covenantal kinship (fellow-Israelite,
  fellow-tribe-member, fellow-Greek as ethnos).  Membership by birth
  into a group larger than the immediate family.
\item[\textbf{E}] Literary / metaphorical extension (personified virtues,
  cosmic-kinship rhetoric, poetic pairings, philosophical universals).
\end{description}

Single-coder classification by the paper's author.  Inter-rater
agreement is not assessed; this is disclosed as a methodological
limitation.

\section*{Results}

\subsection*{Per-bucket proportions with Wilson 95\% CIs}

\begin{center}
\footnotesize
\setlength{\tabcolsep}{4pt}
\begin{tabular}{lrrrrrr}
\toprule
Bucket & $N$ & A \% (95\% CI) & B \% (95\% CI) & C \% & D \% & E \% \\
\midrule
\multicolumn{7}{l}{\textbf{Baseline (non-Christian)}} \\
Papyri (pre-70 CE)     & 35  & 0.0 [0.0, 9.9]    & 100.0 [90.1, 100] & 0.0  & 0.0  & 0.0 \\
LXX                    & 35  & 0.0 [0.0, 9.9]    & 68.6 [52.0, 81.4] & 5.7  & 25.7 & 0.0 \\
Josephus               & 35  & 0.0 [0.0, 9.9]    & 60.0 [43.6, 74.4] & 5.7  & 34.3 & 0.0 \\
Philo                  & 35  & 2.9 [0.5, 14.5]   & 57.1 [40.9, 72.0] & 2.9  & 5.7  & 31.4 \\
Hellenistic literary   & 36  & 0.0 [0.0, 9.6]    & 80.6 [65.0, 90.2] & 0.0  & 2.8  & 16.7 \\
Pre-Pauline literary   & 36  & 2.8 [0.5, 14.2]   & 94.4 [81.9, 98.5] & 0.0  & 0.0  & 2.8 \\
\midrule
\textbf{Baseline aggregate} & \textbf{212} & \textbf{0.9 [0.3, 3.4]} & \textbf{76.9 [70.8, 82.1]} & 2.4 & 11.3 & 8.5 \\
\midrule
\multicolumn{7}{l}{\textbf{Downstream Christian (for context only, not prior)}} \\
NT non-Pauline letters & 35  & 94.3 [81.4, 98.4] & 2.9 [0.5, 14.5]   & 2.9  & 0.0  & 0.0 \\
Gospels + Acts         & 35  & 40.0 [25.6, 56.4] & 51.4 [35.6, 67.0] & 2.9  & 5.7  & 0.0 \\
\midrule
\multicolumn{7}{l}{\textbf{Reference: undisputed Pauline corpus}} \\
Pauline (7 letters)    & 121 & 96.7 [91.8, 98.7] & 0.0 [0.0, 3.1]    & 3.3  & 0.0  & 0.0 \\
\bottomrule
\end{tabular}
\end{center}

\subsection*{Headline finding, with sensitivity analysis}

\begin{center}
\begin{tabular}{lrr}
\toprule
Corpus & A-rate & 95\% CI \\
\midrule
Baseline, sampled registers only ($N=212$)  & 0.9\%  & [0.3\%, 3.4\%] \\
Paul (undisputed, $N=121$)                  & 96.7\% & [91.8\%, 98.7\%] \\
\midrule
\textbf{Nominal ratio (Paul / baseline)}    & \textbf{$\approx 103\times$} & (CIs do not overlap) \\
\bottomrule
\end{tabular}
\end{center}

Read literally, the Pauline A-rate exceeds the sampled non-Christian
baseline by approximately two orders of magnitude.  The nominal ratio
is misleading, however, because the sample excludes the register in
which community-sense \textit{adelphoi} is most heavily attested:
\textbf{Greek voluntary-association inscriptions}.

\subsection*{The inscriptional register and its impact on the baseline}

Greek inscriptions from voluntary cult-associations, mystery-cult
initiate groups, professional guilds, and some civic brotherhoods
across the Hellenistic and Roman East routinely use
\textit{adelphos/adelphoi} as a community-title for fellow members.
This was documented in detail by Philip Harland's 2005 JBL article
\enquote{Familial Dimensions of Group Identity: `Brothers'
(\textit{adelphoi}) in Associations of the Greek East}, and more
broadly in the Kloppenborg-Ascough corpus of Greco-Roman associations
(2011--2020).  A constructed \textit{adelphoi tou theou} / \textit{tou
Serapidos} / \textit{tou Bacchou}-type phrase, parallel in form to
\textit{adelphoi tou kyriou}, is attested in inscriptional material.
This is the single most important piece of context we excluded.

The exclusion was not philosophical but practical: inscriptions are
harder to sample randomly than literary corpora, and the digital
infrastructure for epigraphic corpora is less uniform.  The honest
consequence is that the 0.9\% baseline is a \textit{lower bound}
estimate of the true community-sense rate in Paul's ambient Greek.
Under a plausible correction incorporating inscriptions, the baseline
A-rate would rise into the 5--15\% range (our estimate; we do not
have a rigorous sampled figure).  The Pauline-to-baseline ratio would
correspondingly contract from the nominal $\sim$103$\times$ toward the
$\sim$10--20$\times$ range.  The qualitative conclusion---that Paul's
usage is significantly distinctive against ambient Greek---survives
this correction.  The \textit{quantitative rhetoric of 103$\times$
does not}, and we disavow it here.

\subsection*{What Harland's data changes, and what it strengthens}

Harland's data changes one thing and strengthens another.

\textbf{What it changes.}  The claim that Paul's community-sense
usage is a uniquely Christian semantic innovation is wrong.  Greek
religious associations had been using \textit{adelphoi} for
community-members for centuries before Paul.  Paul's register is
recognisable as a sub-register of Greek religious-associative
discourse; it is distinctive in degree, not in kind.

\textbf{What it strengthens.}  The community-title reading of
Gal~1:19 (\textit{ton adelphon tou kyriou}) and 1~Cor~9:5
(\textit{hoi adelphoi tou kyriou}) gains substantial independent
support from the inscriptional record.  A phrase of the form
\textit{adelphoi tou [patron figure]} designating an inner-circle
community around a focal figure is \textit{attested} in Greek
associations---this is not an exotic reading proposed to save a
thesis; it is the formally normal reading for that specific phrase in
a religious-association register.  The burden of the historicist
reading (biological sibling) is not only that Paul departs from his
own 117-instance pattern, but that he departs from a well-attested
Greek compositional convention to do so.

\section*{What this licenses — and what it doesn't}

\textbf{What it licenses.}  Paul's use of \textit{adelphos} is not
ordinary contemporary Greek usage.  It is a strongly marked dialect
characteristic of early Christian community discourse.  The
\textit{prior} probability that any given Pauline \textit{adelphos}
token is community-sense is very high; the prior probability that any
given Pauline \textit{adelphos} token is biological-sibling-sense is
very low.  This is not an artefact of modern translation or
interpretation; it is a distributional fact of the Pauline corpus
itself, against a clearly-defined non-Pauline non-Christian baseline.

For the disputed \textit{adelphos tou kyriou} passages (Gal 1:19,
1~Cor 9:5), the conservative reading is:
\begin{itemize}[leftmargin=1.5em]
\item In \textit{general} documentary and literary Greek (papyri, LXX,
  Josephus, pagan prose), biological-sibling is overwhelmingly
  dominant; a token picked at random is most likely biological.
\item In the \textit{religious-associative sub-register} (Greek
  voluntary-association inscriptions), community-title \textit{adelphoi}
  with a genitive-of-the-patron-figure is a well-attested
  compositional pattern.
\item Paul's own corpus-wide distribution---117~community, 0~biological
  outside the disputed cases, in 121 tokens---is stable and
  significant, and places his own usage firmly in the
  religious-associative sub-register, not in general Greek.
\item The default reading for \textit{adelphoi tou kyriou} in Paul,
  when all three observations are combined, is the community-title
  reading.  This does not \textit{prove} it; it establishes the
  reading as the parsimonious one and places the interpretive burden
  on the biological-kinship alternative.
\end{itemize}

\textbf{What it does not license.}  This supplement does \textit{not}
compute a posterior probability for the historicity of Jesus.  It does
not multiply base rates by likelihood ratios, because such
multiplication requires auxiliary assumptions (for example, about the
rate at which a community-title in Greek would persist across a
translation from Aramaic, or about the sampling independence of the
$121$ Pauline tokens) that deserve their own methodological defence.
We report the base rates and the deviation; we leave the full
probabilistic argument to readers who wish to construct one.

\section*{Observations on each baseline bucket}

\textbf{Papyri (pre-70 CE).}  100\% biological.  Private letters,
guardianship clauses, co-ownership documents.  The single dominant
pattern in documentary Greek is \textit{adelphos} = biological sibling,
commonly in epistolary greeting-formulas (\enquote{so-and-so to his
brother, greetings}) or legal kinship-designations.  No voluntary-
association uses at all.  The Ptolemaic dating-formula occurrences
(\enquote{Berenice the sister of the gods}) are biological
sister-wives with a cultic epithet, not semantically distinct.

\textbf{LXX.}  68.6\% biological, 25.7\% covenantal.  The D fraction
(\enquote{your brother the Israelite}) reflects Hebrew \textit{'ach}
translated into Greek.  Crucially, no A.  The Jewish scriptural
register has no voluntary-association \textit{adelphos}: the
\enquote{brothers} of the LXX are either patriarchal siblings or
tribal-Israelite compatriots.

\textbf{Josephus.}  60.0\% biological, 34.3\% covenantal.  Similar
profile to LXX, reflecting Josephus's Jewish-historiographical
orientation.  Again, no A.

\textbf{Philo.}  57.1\% biological, 31.4\% metaphorical (virtue-as-
sister, cosmic kinship).  Philo's philosophical register introduces
heavy E (\textit{\enquote{piety and its sister, impiety}}) not seen
elsewhere.  One A (the Therapeutae's brotherhood) is the lone
voluntary-association use across Philo's corpus, and even it is
contested in the critical apparatus.

\textbf{Hellenistic literary (Plutarch, Epictetus, Dio).}  80.6\%
biological, 16.7\% metaphorical.  The E fraction is concentrated in
Epictetus's Stoic cosmopolitan arguments (\enquote{all humans are
brothers, since Zeus is our ancestor}) and Dio's civic-concord
rhetoric.  No A.  The \textit{closest} approximation to community-
sense usage in this bucket is the Stoic \enquote{brothers in the
philosophical life}, classified E because it is a rhetorical
universalism rather than a voluntary membership designation.

\textbf{Pre-Pauline literary (Polybius, Diodorus, Strabo).}  94.4\%
biological.  The one A token in this bucket is striking: Strabo
\textit{Geog.}\ 16.4.21 describes a Nabataean royal court official
called \enquote{brother} (\textit{adelphos}) as an honorific
designation of a favoured companion of the king.  This is the single
closest structural parallel in the entire baseline to a Pauline
\enquote{Brothers of the Lord} title: a voluntary/honorific
\textit{adelphos} designating membership in a recognised inner circle
around a central figure.  It is exactly \textit{one} occurrence in
$212$ baseline tokens.

\textbf{Post-Pauline Christian (for context).}  The NT non-Pauline
letters show the same 94.3\% A distribution as Paul himself, indicating
that the Pauline usage had propagated through early Christian letter-
writing by the late first century.  The Gospels + Acts bucket is
intermediate (40\% A, 51\% B) because narrative frame generates
biographical biological references (Jesus's brothers, Peter-Andrew,
James-John, etc.).  These buckets are \textit{not} prior evidence about
Paul's linguistic environment; they are downstream inheritors of his
(or his community's) distinctive usage.

\section*{Limitations}

\begin{enumerate}[leftmargin=1.5em]
\item \textbf{Single-coder classification.}  All 212 baseline
  classifications were made by the paper's author using the criteria in
  the Pauline supplement.  No inter-rater agreement is assessed.
  A defensible replication would require independent classification by
  at least one other coder and computation of Cohen's $\kappa$.  The
  first-pass tentative classifications from a separate agent (for
  papyri, LXX, Josephus, and the Christian context buckets) have been
  reviewed against the author's coding; disagreements are flagged in
  the data files where they occur.
\item \textbf{Inscriptions excluded --- this is the most important
  limitation.}  Greek inscriptions (IG, SEG, and the specialised corpora
  of voluntary associations edited by Kloppenborg and Ascough) are the
  single register in which community-sense \textit{adelphoi} most
  heavily concentrates.  Harland (JBL 2005, \enquote{Familial Dimensions
  of Group Identity}) documents the inscriptional
  evidence in detail.  Their exclusion means the sampled baseline is a
  \textit{lower bound}: we are confident the true contemporary
  non-Christian A-rate is higher than the 0.9\% sampled figure, and
  likely 5--15\% once inscriptions are included.  This would compress
  the Pauline-to-baseline ratio from the nominal 103$\times$ to roughly
  10--20$\times$.  The qualitative conclusion---Paul's usage is
  distinctive against ambient Greek---survives this correction.  The
  quantitative 103$\times$ figure does not, and should not be cited.
\item \textbf{Sample size modest per bucket.}  $n \approx 35$ per
  bucket gives wide within-bucket CIs.  The aggregate $N = 212$ CI is
  tighter, but expanding to $n = 100$ per bucket would sharpen the
  estimate substantially.
\item \textbf{Bucket-weighting is equal, not register-proportional.}
  An alternative aggregate could weight by estimated volume of Greek
  Paul would have encountered (heavier on papyri and LXX, lighter on
  Plutarch).  Equal weighting across the six buckets gives each
  register equal voice in the aggregate; register-proportional
  weighting would likely push the aggregate A-rate even lower,
  because the register with the largest corpus volume (papyri) has
  0\% A.
\item \textbf{D/E boundary is a judgement call in Philo and Epictetus.}
  Philo's \enquote{sister-virtues} and Epictetus's \enquote{cosmopolitan
  brothers} occupy a fuzzy zone between voluntary-association (A) and
  rhetorical-universalism (E).  We resolved these toward E because the
  referents are abstract (virtues, humanity-at-large) rather than
  concrete members of a named movement.  A stricter coder might push
  some E tokens to A; doing so would raise the baseline A-rate
  modestly but would not change the order-of-magnitude conclusion.
\end{enumerate}

\section*{Conclusion}

In a corpus of $212$ randomly-sampled \textit{adelphos} tokens drawn
from six non-Christian, contemporary-or-earlier Greek corpora (papyri,
LXX, Philo, Josephus, pre-Pauline literary, Hellenistic literary), the
rate of community / voluntary-association usage is $0.9\%$ (95\% CI:
$[0.3\%, 3.4\%]$).  The comparable rate in the seven undisputed Pauline
letters is $96.7\%$ (95\% CI: $[91.8\%, 98.7\%]$).  The nominal ratio
is approximately $103\times$, but this is a lower-bound-of-the-baseline
figure that ignores the Greek inscriptional register
(Harland 2005; Kloppenborg \& Ascough 2011) where community-sense
\textit{adelphoi} concentrates.  Factoring inscriptions into the
baseline would contract the ratio to approximately $10$--$20\times$.

The considered conclusion is twofold.  First, Paul's \textit{adelphos}
usage is distinctive against everyday Greek (where biological meaning
dominates), to a degree that far exceeds sampling noise.  Second, it
is \textit{not} distinctive against the specific sub-register of Greek
religious-associative discourse, where community-title \textit{adelphoi}
with a genitive-of-patron construction (\textit{adelphoi tou theou},
\textit{adelphoi tou Serapidos}) is well-attested.  Paul's usage is
a well-formed token of that sub-register.  Against this combined
picture, any interpretation of a disputed Pauline \textit{adelphos}
token---notably Galatians~1:19---should be anchored in Paul's own
corpus-wide distribution and in the associative-register precedent
for \textit{adelphoi tou X} constructions, not in general-Greek
kinship conventions.  That anchoring is what the main paper's §3.5
argues for on direct lexical grounds; the present supplement
corroborates it quantitatively, with the critical caveat that the
quantitative extremity is an artefact of register-exclusion and should
not be cited in isolation.

\vspace{1.5em}
\noindent\rule{0.5\textwidth}{0.4pt}\par
\vspace{0.5em}
\small\noindent\textit{Attribution:} SBL Greek New Testament
(\textcopyright~2010 SBL/Logos, CC BY 4.0); Perseus canonical-greekLit
and Open Greek and Latin First1KGreek (CC BY-SA 4.0); Duke Databank of
Documentary Papyri (CC BY 3.0).  All source texts and the Python
pipeline that produced the samples are included in the project
repository at \url{https://github.com/sbthomas/historicity}.

\end{document}
